\subsection{THE AUTOMORPHISM GROUP OF A LIE GROUP (ULTRAMETRIC CASE)}

THEOREM 2. When K is ultrametric and locally compact and G is a compact Lie group, assertions (i), (ii), (iii), (v) and (vi) of Theorem 1 are true.

\begin{enumerate}
\def\labelenumi{(\alph{enumi})}
\item
  The uniqueness of the analytic structure considered in (i) is obvious.
\item
  Suppose that G is the Lie group defined by the normable Lie algebra L. Then G is an open and closed ball in L. Let \(w \in \operatorname{Aut} G\) . Then \(\operatorname{L}(w)\) coincides with w in a neighbourhood of 0. Let \(x \in G\) . Let p be the characteristic of the residue field. Then \(p^{n}x\) tends to 0 as n tends to \(+\infty\) . There therefore exists n such that \(w(p^{n}x) = \operatorname{L}(w)(p^{n}x)\) . Therefore
\end{enumerate}

\[
\begin{array}{r l} p ^ {n} w (x) & = w (x) ^ {p ^ {n}} = w (x ^ {p ^ {n}}) = w (p ^ {n} x) \\ & = \mathrm{L} (w) (p ^ {n} x) = p ^ {n} \mathrm{L} (w) (x) \end{array}
\]

whence \(w(x) = \mathrm{L}(w)(x)\) . Thus, \(w = \mathrm{L}(w)|\mathrm{G}\) .

Let \(\Gamma\) be the set of \(\gamma \in \operatorname{Aut} L(G)\) such that \(\gamma(G) = G\) . As \(G\) is open and compact in \(L(G)\) , \(\Gamma\) is an open subgroup of \(\operatorname{Aut} L(G)\) . By the above, Aut \(G\) is identified with \(\Gamma\) , whence there is a Lie group structure on Aut \(G\) , with which properties (i), (ii), (iii) and (vi) of Theorem 1 are obvious. Property (v) follows from Propositions 1 and 3 of no. 1.

\begin{enumerate}
\def\labelenumi{(\alph{enumi})}
\setcounter{enumi}{2}
\tightlist
\item
  We pass to the general case. By § 7, no. 1, Proposition 1, there exists an open compact subgroup \(G_0\) of \(G\) which is of the type considered in (b). Then \(G\) is generated by \(G_0\) and a finite number of elements \(x_1, x_2, \ldots, x_n\) . Let \(\text{Aut}_1G\) be the set of \(u \in \text{Aut}G\) such that \(u(G_0) = G_0\) and \(u(x_iG_0) = x_iG_0\) for \(1 \leqslant i \leqslant n\) . We define as in the proof of Theorem 1, part (c), a semi-direct product \(P\) of \(\text{Aut}G_0\) by \(G_0^n\) and an injective homomorphism \(\zeta\) of \(\text{Aut}_1G\) into \(P\) , whose image is closed in \(P\) .
\end{enumerate}

(d), (e): the argument is exactly as in parts (d), (e) of the proof of Theorem 1 with \(\mathbf{R}\) replaced by \(\mathbf{Q}_p\) and using Proposition 3 instead of Proposition 2.

Remark. If \(\mathbf{K} = \mathbf{Q}_p\) and the Lie group \(\mathbf{G}\) is generated by a compact subset (cf.~Exercise 2), assertions (i), (ii), (iii) and (vi) of Theorem 1 are still true, but not (v) (Exercise 3).

